\begingroup
\setlength{\emergencystretch}{2em}
\capituloTesis{Introducción}
\seccionTesis{El problema}

El manejo del agua forma parte de la conservación de pavimentos porque las variaciones de humedad en la subrasante pueden modificar su respuesta mecánica bajo las cargas del tránsito. Chu et al. (2023) examinaron este proceso mediante estudios de campo, simulaciones de drenaje y ensayos de laboratorio, cuya revisión relacionó las condiciones de humedad con el deterioro inducido por el agua y mostró que la eficacia del drenaje subsuperficial depende de las condiciones en que opera.

La relación entre el drenaje superficial y el estado del pavimento también depende de cómo se evalúan ambos componentes. En sectores urbanos de João Pessoa, Brasil, Diniz y Almeida de Melo (2023) obtuvieron correlaciones débiles entre el índice de condición del pavimento (PCI) y el estado de las cunetas y los sumideros inspeccionados. En Garut, Indonesia, Farida et al. (2025) combinaron la inspección superficial con la modelación hidrológica e hidráulica, que identificó secciones donde los caudales estimados para determinados periodos de retorno superaban la capacidad de los canales.

En el Perú, Mori Linares (2025) documentó cunetas obstruidas por sedimentos y alcantarillas parcialmente colmatadas en la carretera Chamaya–Puente Integración, en Jaén. La evaluación de Salas Cruz (2025) en la carretera Canta–Huayllay se centró en el pavimento, cuya condición caracterizó mediante el PCI, el índice internacional de rugosidad (IRI) y mediciones de deflexión. Ambos estudios ilustran procedimientos de diagnóstico aplicados a componentes que en Simbal se examinarán dentro de las mismas unidades de análisis.

El corredor EMP. LI105 Pedregal–Censo–Cumbray corresponde al proyecto de mejoramiento de la transitabilidad vial interurbana con Código Único de Inversiones 2646655, ubicado en el distrito de Simbal, provincia de Trujillo, departamento de La Libertad (Municipalidad Distrital de Simbal, 2024). La inspección del drenaje distinguirá los recorridos de evacuación longitudinal y los pasos transversales, atendiendo a la conexión entre ambos donde se identifiquen descargas de cunetas hacia alcantarillas de alivio (Ministerio de Transportes y Comunicaciones [MTC], 2014).

En 2026 se evaluará si las diferencias en la condición observable del drenaje se asocian con variaciones del desempeño funcional del pavimento dentro del corredor. Para comparar registros correspondientes a la misma ubicación, se vincularán las observaciones de drenaje con el PCI y el IRI por unidad georreferenciada. El PCI se obtendrá mediante la inspección visual de deterioros superficiales (ASTM International, 2018), mientras que el IRI se calculará a partir del perfil longitudinal para cuantificar la rugosidad (ASTM International, 2021).

\subseccionTesis{Características problemáticas}

\begin{listaTrasTitulo}
\item \textbf{Espacio para la conducción longitudinal.} La inspección deberá distinguir los bordes donde el agua dispone de un recorrido continuo hacia una descarga de aquellos donde el espacio o las interrupciones del borde pueden limitar su conducción (MTC, 2014).
\item \textbf{Abertura de las estructuras transversales.} La condición de cada pase se examinará en sus dos extremos, pues la acumulación de material puede dejar distinta proporción libre en la entrada y en la salida (Mori Linares, 2025).
\item \textbf{Diferencias entre tratamientos del pavimento.} El recapeo sobre una superficie existente y la pavimentación de otros sectores representan antecedentes constructivos que deben reconocerse al interpretar las diferencias de deterioro y rugosidad (Buch et al., 2005).
\item \textbf{Correspondencia espacial entre las mediciones.} El drenaje, el PCI y el IRI se vincularán mediante una delimitación común por sector. La proximidad de las unidades se considerará al contrastar las asociaciones (ASTM International, 2018, 2021; Wagner \& Dray, 2015).
\end{listaTrasTitulo}

\subseccionTesis{Análisis de las características problemáticas}

\apartadoTesis{Espacio para la conducción longitudinal}

La función de una cuneta depende de que pueda captar la escorrentía del borde, conducirla y entregarla a una descarga. El manual del MTC (2014) distingue esas funciones al describir el drenaje longitudinal y prevé alcantarillas de alivio para evacuar el agua recogida por las cunetas. Si el espacio lateral se reduce, la inspección debe reconocer por dónde puede continuar el agua y dónde termina el recorrido, ya que la presencia de un segmento de cuneta aporta información parcial sobre su conexión con la descarga. En Simbal interesa establecer qué trayectorias pueden seguirse en cada margen, especialmente donde se encuentre un borde estrecho, un cambio de sección o una terminación cerrada.

La observación se trasladará a una medida de continuidad libre mediante el cociente entre la longitud continua y la longitud evaluada. Para sostener la comparación entre unidades, la ficha registrará las longitudes, la posición de las interrupciones y el punto de descarga cuando sea identificable. La continuidad se interpretará como condición geométrica observable en la fecha de inspección. La estimación de capacidad hidráulica exigiría incorporar, entre otros parámetros, la pendiente, la sección y la rugosidad del recorrido, tal como prevé el método de diseño de cunetas del MTC (2014). Esta distinción permite examinar la relación con PCI e IRI sin atribuir un caudal a cada segmento a partir de su longitud.

\apartadoTesis{Abertura de las estructuras transversales}

En una estructura transversal, los sedimentos o residuos pueden ocupar parte del espacio de paso en cualquiera de sus extremos. El diagnóstico de Mori Linares (2025) en Jaén documentó alcantarillas parcialmente colmatadas, mientras el manual de hidrología y drenaje vincula el diseño de la obra transversal con el paso del flujo y del material que transporta (MTC, 2014). La inspección del pase en Simbal abarcará la entrada y la salida, de modo que pueda distinguirse una reducción localizada en un extremo de una afectación visible en ambos. La identificación del extremo más restringido orientará la descripción de la estructura y el registro fotográfico.

El indicador previsto expresará la menor proporción de abertura libre entre los dos extremos. Para obtenerla se registrará la sección interior de cada abertura y la superficie ocupada, recurriendo a fotografías frontales con escala cuando la geometría dificulte la medición directa. La coincidencia entre fotografía y ficha permitirá reconstruir el cálculo, identificar el extremo que determina el indicador y conservar las observaciones complementarias sobre erosión o integridad. La lectura de esa proporción se mantendrá en el terreno de la inspección física. Un análisis hidráulico del paso requeriría examinar también el conducto, los niveles y las condiciones de descarga, de acuerdo con los criterios del MTC (2014).

\apartadoTesis{Diferencias entre tratamientos del pavimento}

El recapeo se ejecuta sobre una superficie existente, cuyo perfil y preparación pueden incidir en la regularidad de la capa colocada. A partir de datos del programa estadounidense de desempeño de pavimentos a largo plazo, Buch et al. (2005) estudiaron cómo la condición previa y las características del tratamiento se relacionan con la rugosidad posterior al recapeo. En Simbal, el antecedente constructivo aporta contexto para interpretar las diferencias entre sectores sometidos a tratamientos distintos. Una diferencia de IRI entre dos sectores puede coincidir con distinta continuidad de drenaje y con una historia constructiva diferente, situaciones que deben reconocerse al examinar la distribución de las observaciones.

La identificación del tratamiento por sector acompañará los registros de PCI e IRI sin alterar sus procedimientos de cálculo. Los deterioros se registrarán mediante la inspección prevista para el PCI, mientras que la rugosidad se obtendrá del perfil longitudinal conforme a la práctica correspondiente (ASTM International, 2018, 2021). Cuando una unidad se encuentre próxima a una transición constructiva, la ubicación del cambio será relevante para la lectura de los indicadores agregados. El antecedente de intervención servirá para contextualizar las asociaciones, dentro de un diseño orientado a comparar condiciones observadas en una campaña.

\apartadoTesis{Correspondencia espacial entre las mediciones}

El drenaje transversal se observa en la estructura, la evacuación longitudinal se mide en los bordes y las propiedades del pavimento se registran en la superficie vial. La correspondencia entre esas observaciones depende de que cada ficha y cada perfil puedan asignarse al mismo sector, puesto que el PCI se obtiene sobre una superficie inspeccionada y el IRI sobre una trayectoria longitudinal (ASTM International, 2018, 2021). El pase de agua servirá de referencia para delimitar la unidad estructura–sector, cuyo código vinculará las observaciones de ambos márgenes, la abertura de los extremos y los indicadores funcionales.

En estructuras próximas, el criterio de delimitación por progresivas evitará asignar la misma superficie a dos unidades. Esa separación física permite identificar el soporte de las mediciones, aunque los sectores vecinos pueden compartir características del relieve, del pavimento o del drenaje. Wagner y Dray (2015) muestran cómo construir distribuciones nulas que incorporen la estructura espacial en contrastes con observaciones georreferenciadas. En este estudio, la posición de cada unidad permitirá considerar la proximidad en los contrastes previstos y distinguir la magnitud descriptiva de las asociaciones de la evidencia estadística obtenida al reconocer la estructura espacial.

\seccionTesis{Enunciado del problema}

¿Qué relación existe entre la condición del drenaje vial y el desempeño funcional del pavimento en el tramo EMP. LI105 Pedregal–Censo–Cumbray, Simbal, durante 2026?

\seccionTesis{Alcance}

La investigación comprende el corredor EMP. LI105 Pedregal–Censo–Cumbray, en el distrito de Simbal, provincia de Trujillo, identificado como parte del proyecto de mejoramiento de la transitabilidad vial interurbana con Código Único de Inversiones 2646655 (Municipalidad Distrital de Simbal, 2024). El análisis se circunscribe a las observaciones previstas para 2026 y se organiza alrededor de las estructuras de drenaje transversal y de los sectores de pavimento asociados a ellas. La unidad estructura–sector permite reunir en una misma observación el pase de agua, los recorridos longitudinales y la superficie evaluada, conforme a límites establecidos por tramo y progresiva.

La condición del drenaje se examinará mediante la continuidad longitudinal libre y la proporción de abertura libre en los extremos del pase. Las fichas de campo conservarán las longitudes, las áreas, la ubicación y las fotografías que permitan revisar cada indicador. El alcance de estas observaciones comprende el estado físico visible de los recorridos y de las aberturas. Quedan fuera la estimación de caudales de diseño, la simulación de crecidas y la evaluación de la capacidad hidráulica para periodos de retorno, que requieren los parámetros y procedimientos propios del diseño de drenaje (MTC, 2014).

El desempeño funcional se caracterizará mediante la condición superficial y la rugosidad longitudinal. Las fichas de deterioros aportarán los registros para calcular el PCI, mientras que los perfiles de las huellas proporcionarán la información necesaria para calcular el IRI (ASTM International, 2018, 2021). La evaluación se circunscribe a esas dos propiedades de la superficie. La capacidad estructural, la resistencia al deslizamiento y la evolución del deterioro en campañas sucesivas quedan fuera del alcance establecido.

El diseño transversal y correlacional delimita los contrastes a las cuatro combinaciones entre los indicadores de drenaje y los de pavimento. El análisis estimará la magnitud y el sentido de las asociaciones entre sectores, con atención a la proximidad espacial en el contraste de hipótesis (Wagner \& Dray, 2015). Los resultados corresponderán a las unidades examinadas y a las condiciones registradas durante el periodo del estudio. Su interpretación reconocerá el antecedente constructivo de los sectores y las condiciones observadas en la campaña, sin extender las estimaciones a otros corredores ni establecer efectos causales o tasas de deterioro.

\seccionTesis{Viabilidad}

La ejecución del estudio se organiza en una inspección del drenaje, una evaluación visual del pavimento y un perfilado longitudinal sobre los sectores delimitados. Las fichas, la cinta de medición y las fotografías con escala permiten documentar longitudes y aberturas con recursos de campo de uso directo. La evaluación del pavimento dispone de procedimientos definidos para el inventario visual de deterioros y el procesamiento del perfil longitudinal (ASTM International, 2018, 2021). La disponibilidad del equipo de perfilado y la comprobación de su respuesta antes de medir constituyen condiciones concretas de la campaña.

El soporte documental parte de la identificación pública del corredor y de la intervención vial (Municipalidad Distrital de Simbal, 2024). La preparación requiere contrastar el inventario de estructuras con las progresivas del recorrido y verificar en campo el acceso a los extremos de cada pase. Las mediciones y fotografías generarán registros vinculados al código de unidad, cuya revisión permitirá detectar inconsistencias de ubicación, dimensiones o identificación antes del análisis. La disponibilidad de documentos técnicos y la confirmación del acceso condicionan el grado de detalle con que podrá describirse el antecedente constructivo de cada sector.

La organización de las jornadas debe contemplar el tránsito, la accesibilidad de las estructuras y las condiciones en que puedan realizarse lecturas comparables. La inspección se programará con las coordinaciones y permisos que correspondan para ocupar temporalmente los puntos de medición. Las fichas y fotografías se limitarán a la información necesaria para identificar la infraestructura y respaldar las mediciones. La previsión de una revisión de fichas y una prueba piloto permite comprobar anticipadamente la claridad de los criterios de inspección, la repetición de mediciones y la correspondencia de los registros.

La secuencia prevista para 2026 separa preparación, prueba piloto, levantamiento, procesamiento y análisis, de manera que una discrepancia de campo pueda revisarse antes de consolidar los indicadores. La ejecución económica depende de la disponibilidad del equipo, los traslados y el tiempo de los evaluadores para completar las inspecciones y los perfiles. El procesamiento podrá comenzar una vez que las observaciones de drenaje, las fichas PCI y los perfiles correspondientes a cada unidad estén completos y conciliados. Con esa estructura de datos se podrán examinar los cuatro pares de indicadores y aplicar los contrastes previstos con incorporación de la proximidad entre unidades (Wagner \& Dray, 2015).

\seccionTesis{Justificación}

La conservación del corredor requiere información que permita ubicar las restricciones de evacuación y examinar el estado de la superficie en los sectores correspondientes. El manual del MTC (2014) explica la función de los recorridos de captación, conducción y descarga, mientras los procedimientos de evaluación de pavimentos distinguen entre deterioros visibles y rugosidad longitudinal (ASTM International, 2018, 2021). La pregunta de investigación surge al reunir esas observaciones en una misma unidad espacial y examinar cuánto varían conjuntamente en el corredor.

La asociación encontrada entre el drenaje superficial y la condición del pavimento en otros entornos no establece qué magnitud tendrá esa relación en Simbal. Diniz y Almeida de Melo (2023) documentaron asociaciones débiles en los sectores urbanos que inspeccionaron, de modo que su estudio fundamenta la pertinencia de medir ambas condiciones, con procedimientos acordes al corredor analizado. El registro de continuidad longitudinal y abertura transversal permitirá examinar por separado cada componente en relación con el PCI y el IRI.

El interés de la evaluación comprende la oportunidad de disponer de un diagnóstico localizado, su posible aporte a la conservación de una vía interurbana y la evidencia que se obtenga sobre las relaciones entre drenaje y pavimento. Al documentar cómo se asigna cada observación al sector correspondiente, el estudio establecerá una base común para revisar las mediciones, interpretar el alcance de las asociaciones y valorar su posible uso en la programación de inspecciones. La utilidad de esta información dependerá de la cobertura de los registros y de la consistencia de los patrones que se obtengan.

\subseccionTesis{Conveniencia}

El corredor Pedregal–Censo–Cumbray constituye un caso delimitado por la intervención de mejoramiento de la transitabilidad vial de Simbal (Municipalidad Distrital de Simbal, 2024). La localización de los pases de agua proporciona referencias para organizar la inspección en sectores identificables. Dentro de cada unidad será posible consultar la ficha de drenaje junto con la evaluación del pavimento, lo que facilitará la revisión de coincidencias y discrepancias entre ambas condiciones durante la campaña de 2026.

Una ficha de continuidad o de sección libre aporta una observación que puede verificarse en el sitio documentado. La utilidad de combinarla con el PCI y el IRI consiste en comparar esa observación con dos lecturas distintas de la superficie, una derivada de los deterioros visibles y otra del perfil longitudinal (ASTM International, 2018, 2021). La codificación común permitirá que los responsables de revisar el estudio puedan volver a la evidencia de un sector cuando sus indicadores requieran una explicación contextual. Esta organización también permitirá identificar omisiones de campo antes de atribuir una diferencia a las condiciones del lugar.

La oportunidad del estudio se sitúa en disponer de una caracterización del corredor para un periodo definido, susceptible de consulta posterior. Los profesionales que intervienen en su conservación podrán utilizar el inventario sectorial para seleccionar puntos de reinspección y examinar si las observaciones de drenaje coinciden con deterioro o rugosidad en la superficie asociada. El manual del MTC (2016) contempla actividades diferenciadas de conservación del drenaje, entre ellas la limpieza de cunetas y alcantarillas. El estudio aportará información para examinar dónde corresponde verificar esas necesidades, mientras la elección de una intervención dependerá de la inspección específica y de las condiciones que encuentre el responsable de la vía.

\subseccionTesis{Relevancia social}

La intervención priorizada en Simbal se orienta a mejorar la transitabilidad del corredor que vincula Pedregal, Censo y Cumbray (Municipalidad Distrital de Simbal, 2024). La información que produzca el estudio se relaciona con ese servicio por su aplicación potencial a la conservación de la vía. Entre sus usuarios se encuentran los residentes que se desplazan entre los sectores y los operadores que recorren el corredor, mientras que la entidad responsable y su personal técnico constituyen los destinatarios inmediatos de las fichas y mediciones.

La rugosidad longitudinal describe una propiedad de la superficie vinculada con la respuesta de un vehículo al circular sobre un perfil, según el procedimiento que sustenta el IRI (ASTM International, 2021). El PCI, por su parte, permite localizar y valorar deterioros visibles dentro de la superficie inspeccionada (ASTM International, 2018). Al relacionar ambas mediciones con las observaciones de cunetas, recorridos y pases de agua, el diagnóstico podrá señalar en qué sectores conviene comprobar conjuntamente el estado de la rodadura y la evacuación, en lugar de orientar la revisión a partir de un promedio del corredor.

El beneficio directo del trabajo será la disponibilidad de registros localizados para la inspección técnica. Un beneficio indirecto podría surgir si esa información se utiliza para verificar puntos de acumulación de agua, programar labores de conservación o revisar sectores con deterioro o rugosidad que merezcan seguimiento. La limpieza de cunetas y alcantarillas figura entre las actividades de mantenimiento rutinario descritas por el MTC (2016), lo cual permite vincular las observaciones con posibles acciones de comprobación en el terreno. Los cambios en tiempos de viaje, seguridad, acceso a servicios o costos de transporte requerirían mediciones adicionales. La relevancia social se fundamenta en la posibilidad de orientar decisiones sobre una infraestructura de uso interurbano mediante observaciones que pueden localizarse y revisarse.

\subseccionTesis{Justificación teórica}

El aporte teórico consiste en examinar la relación entre dos componentes que suelen describirse con medidas distintas. La literatura sobre drenaje subsuperficial explica la interacción entre humedad y respuesta del pavimento, mediante evidencia procedente de estudios de campo, simulación y laboratorio (Chu et al., 2023). Esa explicación permite identificar mecanismos pertinentes para la conservación, mientras que la medición prevista en Simbal se refiere a la condición física observable de los recorridos y las aberturas del drenaje superficial. El estudio relacionará esa condición con el deterioro visible y la rugosidad en los sectores evaluados.

La evidencia de Diniz y Almeida de Melo (2023) aporta un precedente de comparación entre drenaje superficial y PCI en vías urbanas. Farida et al. (2025) trabajaron desde otra perspectiva al incorporar modelación hidráulica en el diagnóstico de canales y pavimento. La investigación en Simbal examinará las asociaciones a partir de mediciones geométricas de continuidad y abertura libre, con registros del PCI y del IRI correspondientes a una misma unidad. Ese planteamiento permitirá reconocer si los patrones entre drenaje y condición superficial coinciden con los observados para la rugosidad o si las dos propiedades de la superficie presentan comportamientos distintos dentro del caso.

El aporte será evidencia empírica situada sobre la magnitud y la dirección de cada relación. La distinción entre dimensiones ayudará a interpretar resultados que podrían perderse si la evaluación se redujera a una calificación general de la vía. El indicador de rugosidad y el de condición superficial provienen de procedimientos diferenciados (ASTM International, 2018, 2021), por lo que sus asociaciones se discutirán atendiendo a qué mide cada uno y a las condiciones constructivas y espaciales que coincidan en los sectores.

\subseccionTesis{Justificación metodológica}

El aporte metodológico se concentra en la correspondencia espacial entre observaciones obtenidas mediante instrumentos distintos. La ficha de drenaje longitudinal registrará la longitud libre respecto de la longitud examinada y la ficha transversal documentará la abertura disponible en cada extremo del pase. Los formatos de inspección y los perfiles longitudinales permitirán calcular de manera independiente el PCI y el IRI (ASTM International, 2018, 2021). La clave de unidad, formada a partir del tramo y la progresiva de referencia, permitirá relacionar cada indicador con los registros que lo sustentan.

La delimitación de sectores exclusivos mantiene identificables las superficies y trayectorias incorporadas en cada observación. Para examinar la consistencia del registro, la metodología prevé revisión de la ficha de drenaje y mediciones independientes en unidades seleccionadas, junto con comprobaciones de repetición del perfilado. El respaldo de fotografías y archivos de medición permitirá revisar diferencias entre evaluadores o recorridos antes de consolidar los indicadores. Las reglas podrán adaptarse a otros corredores con estructuras identificables y observaciones de pavimento que puedan asignarse con precisión al sector asociado.

La comparación entre sectores también requiere considerar que la cercanía puede introducir dependencia entre observaciones. La aleatorización espectral de Moran permite construir distribuciones nulas que incorporan la estructura espacial en contrastes de asociación (Wagner \& Dray, 2015). Su aplicación prevista complementará la descripción mediante coeficientes de rangos y permitirá interpretar por separado la magnitud observada y el respaldo estadístico de cada hipótesis. La contribución reside en documentar la secuencia que conecta medición, delimitación espacial y contraste, de modo que pueda revisarse cómo se obtuvo cada dato y qué condiciones se tomaron en cuenta al interpretarlo.

\subseccionTesis{Justificación práctica}

El diagnóstico podrá orientar la programación de inspecciones del corredor al identificar, por tramo y progresiva, los puntos en los que se registren interrupciones del recorrido longitudinal, reducciones de abertura o diferencias en la condición del pavimento. El MTC (2016) contempla la limpieza de cunetas y alcantarillas como actividades de mantenimiento rutinario. Su aplicación requiere examinar las condiciones de cada elemento, de modo que las fichas y fotografías previstas puedan servir como insumos para comprobar si el hallazgo corresponde a material acumulado, a una discontinuidad del recorrido o a otra condición física observable.

La lectura conjunta de los indicadores permitirá decidir qué elementos conviene revisar en un mismo sector. Si una abertura reducida coincide con deterioros de la superficie, la inspección posterior podrá verificar tanto los extremos del pase como el tipo y la extensión del daño. Si la observación principal corresponde a la rugosidad, la consulta del perfil y del antecedente constructivo ayudará a precisar dónde concentrar esa revisión. Estos usos se apoyan en la finalidad del PCI para orientar prioridades de mantenimiento y del IRI para aportar información sobre la regularidad de la rodadura (ASTM International, 2018, 2021).

La entidad responsable y el personal técnico de conservación podrán utilizar el inventario como referencia para comparar inspecciones, justificar la selección de sectores y documentar cambios en visitas posteriores. La programación de intervenciones y su dimensionamiento exigirán información de ingeniería adicional cuando la condición observada exceda el mantenimiento rutinario. En otros corredores con pases de agua identificables, las fichas y reglas de correspondencia espacial podrán servir como referencia para organizar un levantamiento comparable, previa adaptación a la configuración del drenaje y a las características del pavimento local.
\endgroup
